\section{Introduction}

Marginum measures the fluorescence of tissue pieces while they are aspirated during surgery.
A dye makes tissue fluoresce, and the signal is recorded by two sensors. Each tissue piece
that passes the sensor gives a peak in the signal. % CHECK: confirm the clinical background with the course material.
Knowing how many pieces were aspirated, and how much tissue they make up, is useful because
the surgeon cannot see this during the operation. % CHECK: confirm this motivation.

The recordings are CSV files with two integer columns, \texttt{adc1} and \texttt{adc2},
sampled at \SI{50}{\kilo\hertz}. A single recording can last several hours and be
\SI{7.6}{\giga\byte} large. This makes the data hard to handle in three ways:
\begin{itemize}
  \item The files are too large to open in a normal way, and the result is needed quickly.
  \item The signal contains noise and a slowly drifting baseline, so the peaks cannot be
        read from the raw values.
  \item A peak does not tell by itself what it comes from. It can be dyed tissue, tissue without
        dye, or water, and one piece often gives several peaks.
\end{itemize}

The goal of this project is to find each aspirated tissue piece, decide whether it is tissue
(with or without dye) or water, and estimate the amount of tissue. The work is divided into
two parts, which follow each other:
\begin{description}
  \item[Part A] reads and cleans the signal. It reduces the recording to
        \SI{1}{\kilo\hertz}, removes noise and baseline drift, and measures the noise level
        $\sigma$ from which the detection threshold is set (Section~\ref{sec:threshold}).
  \item[Part B] uses the cleaned signal and $\sigma$ to detect the peaks, group them into pieces,
        label them as dyed tissue, undyed tissue or water
        (Section~\ref{sec:detection-methods}), and estimate the amount of tissue
        (Section~\ref{sec:amount-methods}).
\end{description}
An interactive Streamlit application lets the user browse a recording and inspect every
detected piece.

The report follows the course outline. Section~\ref{sec:recordings} describes the recordings,
and the methods, results and discussion are given for each part in turn. We used 29 recordings,
some with a known answer (the type of every piece, or the amount aspirated), so the methods
could be checked against the truth.